\documentclass[10pt,a4paper]{amsart}
\usepackage[margin=19mm]{geometry}
\usepackage{amsmath,amssymb,mathtools}
\usepackage[scr=rsfs,cal=euler]{mathalfa}
\usepackage{xfrac}
\usepackage[unicode,hidelinks,bookmarks=false]{hyperref}
\newcommand{\ie}{i.e.\ }
\newcommand{\cf}{cf.\ }
\newcommand{\resp}{respectively\ }
\newcommand{\Subsection}{\subsection}
\providecommand{\nobreakdash}{\nobreak\hbox{-}}
\setlength{\emergencystretch}{2em}
\multlinegap0pt
\usepackage{amssymb}
\usepackage{mathtools}
\usepackage[shortlabels]{enumitem}
\usepackage{xcolor}
\newtheorem{thm}{Theorem}
\newcommand{\Rc}{\mathrm{Rc}}
\newcommand{\Rm}{\mathrm{Rm}}
\newcommand{\dV}{\mathrm{d}V}
\title{A note on Ricci flow from small curvature concentration and a Morrey-type condition}
\author{Albert Chau, Adam Martens}
\date{Source: arXiv:2603.28976v1 (CC BY 4.0)}
\begin{document}
\maketitle

\noindent\emph{Source and licence.} A note on Ricci flow from small curvature concentration and a Morrey-type condition. Version arXiv:2603.28976v1 (CC BY 4.0), distributed by arXiv under the Creative Commons Attribution 4.0 International licence, http://creativecommons.org/licenses/by/4.0/, as recorded in the arXiv licence metadata for this exact version and shown on the abstract page https://arxiv.org/abs/2603.28976v1. Full licence notice: \texttt{licenses/arxiv-2603.28976-CC-BY-4.0.txt}.\par\smallskip

\noindent\textit{Editorial scope.} Counted unit: the article's thm upperContLp (source0.tex lines 417--427). The article's complete original proof of it is delivered with it (source0.tex lines 429--527, one \texttt{proof} environment of 99 lines). No mathematics is written, repaired or re-derived: the statement and the proof are the article's own lines, in the article's own notation, and no retained line is substituted or re-flowed.  Retained uncounted as the source-local context the proof needs: lines 191--198.  Nothing was omitted from any retained span, so the package carries no omission record.  The retained lines cite 3 works, whose entries are transcribed verbatim from the article's own bibliography source in the e-print and delivered with short editorial labels recorded in the item's reference mapping.  No retained line contains a figure environment, an image-inclusion command or drawing code, so no image is delivered and the package declares no asset dependency.  Editorial formatting only, in addition: an AMS \texttt{amsart} preamble that restates the article's own declarations and macros used by the retained lines, namely the environments \texttt{thm} on separate counters, together with the standard packages those lines need.  The retained environments print in this document as Theorem 1 in order of appearance, whereas the article prints the same items with the numbering of its own full document.  The complete native source of the article as distributed in the arXiv e-print for this exact version is bundled unchanged as \texttt{sources/197481/source0.tex}.
\begin{equation}\label{RF}
                \left\{
                \begin{aligned}
                        \frac{{\partial }}{{\partial t}}g(t) &= -2 \text{Rc}_{g(t)}, \\
                        g(0) &= g
                \end{aligned}
                \right.
\end{equation}

    \begin{thm}\label{upperContLp}
        For any $n\geq 3$, $C_0>0$ and $a\in (0,1)$, there exists $\widetilde T(n,a,C_0)>0$ such that the following holds. Let $(M,g(t))$, $t\in [0,T]$ be any complete Ricci flow (not assumed bounded curvature), which satisfies the integrable to time-0 curvature bound
        \begin{equation}\label{integrableCurveBoundLemma}
        \|\Rm\|_{g(t)}\leq \frac{C_0}{t^{1-a}}
        \end{equation}
        for all $t\in (0,T]$. Then the global $L^p$-curvature ($p\geq 1$) of the solution $g(t)$ satisfies
        \[
        \int_M |\mathrm{Rm}|_{g(t)}^p\,\dV_{g(t)}\leq e^{\frac{C_0(16p+n)}{a}t^a}\int_M |\mathrm{Rm}|_{g(0)}^p\,\dV_{g(0)}
        \]
        for all $t\in [0,\min\{T,\widetilde T\}]$, provided the RHS above is finite.
    \end{thm}

\begin{proof}
For clarity of proof, we begin by defining 
\[
\widetilde T=\min\{(nC_0)^{-1/a},\widehat T(n,1,3)\},
\]
where $\widehat T$ is from {\cite[Theorem 1.1]{LeeTam}}. Also for brevity, write $g=g(0)$. Fix some $p\geq 1$ and without loss of generality, assume that $|\Rm|_{g(0)}^p$ is integrable on $(M^n, g)$. Let $v_0 : M \to \mathbb{R}$ be an integrable function on $(M^n, g)$ satisfying $0 < v_0 \le 1$. For any $\varepsilon \in (0,1]$, define $v_{\varepsilon}(x, t)$ to be the solution to the initial value problem
\[
\begin{cases}
\partial_t v_{\varepsilon} = \Delta_{g(t)} v_{\varepsilon},\\
v_{\varepsilon}(x, 0) = \varepsilon v_0(x).
\end{cases}
\]
Also define
\[
K_{\varepsilon}^2 := | \mathrm{Rm} |_{g(t)}^2 + v_{\varepsilon}^2.
\]


In {\cite[Lemma 2.1]{ChauMartens}} the following estimate was shown using the standard estimate for the evolution of the curvature tensor along Ricci flow:
\[
\partial_t K_{\varepsilon}^p \le \Delta_{g(t)} K_{\varepsilon}^p + 8 p K_{\varepsilon}^{p+1},
\]
which, after applying the curvature bound in \eqref{integrableCurveBoundLemma}, we can simplify to
\begin{equation}\label{evolOfKp}
\partial_t K_{\varepsilon}^p \le \Delta_{g(t)} K_{\varepsilon}^p + 16C_0p t^{-1+a}\, K_{\varepsilon}^{p}.
\end{equation}
Note, the addition of $v_{\varepsilon}^2$ above is only needed to prevent possible division by zero in deriving the above estimates in the case when $p<2$.  Thus if we define $\widetilde K_{\varepsilon}^p:=e^{-(16C_0p/a)t^{a}}K_{\varepsilon}^p$, then $\widetilde K_\epsilon^p$ is a subsolution to the time dependent heat equation:  

\begin{equation}\label{upperlowerK}
\partial \widetilde K_{\varepsilon}^p  \le  \Delta_{g(t)} \widetilde K_{i,\varepsilon}^p \text{   on } M\times[0, T].
\end{equation}

    

Now we fix a sequence of nested compact sets $\{\Omega_j\}_{j=0}^\infty$ exhausting $M$, and for each $j$, let $\phi_j: M \to [0,1]$ be a smooth cutoff function  supported in $\Omega_j$ and with $\phi_j \equiv 1$ in $\Omega_{j-1}$. For any $\varepsilon > 0$ and $j$ we denote by 
\[
u_{j} : \Omega_j \times [0,T] \to \mathbb{R}
\] 
the solution to the initial value Dirichlet boundary value problem
\[
\begin{cases}
\partial_t u_{j} = \Delta_{g(t)} u_{j}, & \text{on }  \Omega_t\times[0, T] \\
u_{j}(x,0) = \phi_j K_{\varepsilon}^p(x,0), & x \in \Omega_j\\
u_{j}(x,t) = 0, & x \in \partial \Omega_j.
\end{cases}
\]
Here the dependence of $u_j$ on $\varepsilon$ has been suppressed for simplicity and will not cause confusion for our purpose.  Just as in the proof of Proposition 2.2 in \cite{ChauMartens}, we obtain the evolution inequality
\[
\frac{\mathrm{d}}{\mathrm{d}t}\int_{\Omega_j} u_{j} \,\dV_{g(t)} \le n\int_{\Omega_j} u_{j}|\mathrm{Rm}|_{g(t)} \,\dV_{g(t)}\le nC_0 t^{-1+a} \int_{\Omega_j} u_{j} \,\dV_{g(t)},
\]
where we have used \eqref{integrableCurveBoundLemma} and the scalar curvature bound $|\mathrm{Scal}|\leq n|\Rm|$. Thus by comparison to the corresponding ODE, we have
\[
\int_{\Omega_j} u_{j}(x, t) \,\dV_{g(t)}\leq e^{(nC_0/a)t^{a}}\int_{\Omega_j} u_{j}(x, 0) \,\dV_{g(0)} \le e^{(nC_0/a)t^{a}}\int_{M} |\mathrm{Rm}|_{g(0)}^p  \,\dV_{g(0)} .
\]

By the classical maximum principle, the sequence $u_j (x, t)$ is non-decreasing in $j$ for any fixed $(x, t)$.  On the other hand, the sequence $u_j(x, t)$ must be bounded above on any compact set of $M$ independent of $j, t$ as otherwise the integral upper bound above would be violated by the classical Harnack inequality for linear parabolic equations in \cite{KrylovSafanov}.  It follows that $u_j(x, t)$ converges as $j\to \infty$ in $L^{\infty}_{loc}(M\times[0, T])$ to a limit $u(x, t)$, and by the parabolic Schauder estimates it follows that $u(x, t)\in C^{\infty}$ and solves

\[
\begin{cases}
\partial_t u= \Delta_{g(t)} u\; \text{ on }M\times[0, T] & \\
u_{j}(x,0) = K_{\varepsilon}^p(x,0), \text{ on }M,& \\
\end{cases}
\]
while satisfying
\begin{equation}\label{integrableu}
\int_{M} u(x, t) \,\dV_{g(t)} \le e^{(nC_0/a)t^{a}}\int_{M} K_{\varepsilon}^p(x,0)  \,\dV_{g(0)} .
\end{equation}
From the above equation for $u$ and \eqref{upperlowerK} we get that the function $\varphi^R(x,s):=\widetilde K_{\varepsilon}^p(x,Rs)-u(x,Rs)$ satisfies
\begin{equation}\label{heatphiR}
\begin{cases}
\partial_s \varphi^R\leq \Delta_{g_R(t)} \varphi^R \; \text{ on }M\times[0, T/R] & \\
\varphi^R(x, 0)=0, \text{ on }M& \\
\end{cases}
\end{equation}
for any $R>0$ where $g_R(x, t):=g(x, Rt)/R$ is also a solution to the Ricci $\eqref{RF}$ (with initial condition $g(0)/R$).
\,

For simplicity of notation, we will – for the remainder of the proof – assume that $T\leq \widetilde T$. Note that the integral bound \eqref{integrableCurveBoundLemma}, together with the bound $\widehat T\leq (nC_0)^{-1/a}$, implies the scale-invariant curvature bound
\[
\|\Rm\|_{g(t)}\leq \frac{1}{nt} \quad \implies \quad \Rc_{g(t)}\leq \frac{g(t)}{t}
\]
on $(0,T]$.  Now recall that we are assuming $T\leq \widetilde{T}\leq \widehat T$ where $\widehat T = \widehat T(n, 1, 3)$ is from the conclusion of the local maximum principle in {\cite[Theorem 1.1]{LeeTam}} with $\alpha, l$ there being $1,3$ respectively.  From that theorem, it follows that whenever $R\geq 1$, the statement of that theorem gives at any point $p\in M$ the inequality
\[
\varphi^R(p,s)\leq s^3 \,\,\, \text{for all} \,\,\, s \in [0, R^{-1}T].
\]
Fixing any $t_0\in [0, T]$ and evaluating the above at $s_R:=t_0/R$ gives 
\[
\widetilde K_{\varepsilon}^p(p,t_0)-u(p,t_0)\leq t_0^3 R^{-3}. 
\]
Letting $R\to \infty$ and noting $(p, t_0)$ was chosen arbitrarily, we conclude that
\begin{equation}\label{upperboundu}\widetilde K_{\varepsilon}^p-u\leq 0 \,\,\, \text{on} \,\,\, M\times[0, T].
\end{equation}
\,

Thus, using \eqref{integrableu} we conclude 

\[\int_{M} \widetilde K_{\varepsilon}^p(x, t) \,\dV_{g(t)} \le e^{(nC_0/a)t^{a}}\int_{M} K_{\varepsilon}^p(x,0)  \,\dV_{g(0)} \]
for all $t\in [0, T]$. The proof is then readily completed by recalling $\widetilde K_{\varepsilon}^p = e^{-(16C_0 p/a)t^a)}K_{\varepsilon}^p$ and  taking $\epsilon\to 0$. \\
\end{proof}

\begin{thebibliography}{99}
\bibitem[ChauMa]{ChauMartens} A.~Chau and A.~Martens, \textit{Long-time Ricci flow existence and topological rigidity from manifolds with pinched scale-invariant integral curvature}, to appear in Comm. Anal. Geom., arXiv:2403.02564.
\bibitem[Krylov]{KrylovSafanov} N.~V.~Krylov and M.~V.~Safonov, \textit{A property of the solutions of parabolic equations with measurable coefficients}, Izv. Akad. Nauk SSSR Ser. Mat. \textbf{44} (1980), no.~1, 161--175.
\bibitem[LeeTam]{LeeTam} M.-C.~Lee and L.-F.~Tam, \textit{Some local maximum principles along Ricci flows}, Canad. J. Math. \textbf{74} (2022), no.~2, 329--348.
\end{thebibliography}
\end{document}
